\section{两条技术路线：In-Context 与 End-to-End}

上一章定义了 Agent 和环境类型，本章讨论怎样把 Agent 做出来。节目把路线分为 In-Context Learning 和 End-to-End Training。In-context 路线依赖强预训练模型、prompt、工具协议和上下文工程，迭代快、灵活性高；End-to-end 路线把行为写进模型权重，推理更稳定，但训练成本高、环境构建复杂。

\begin{figure}[H]
\centering
\includegraphics[width=0.94\textwidth]{agent-training-routes.png}
\caption{Agent 两条训练路线：In-context 灵活，End-to-end 稳定但成本高。自制概念图，依据 00:14:50--00:30:57 对谈内容整理。}
\end{figure}

\begin{knowledgebox}{读图：两条路线不是非此即彼}
很多真实系统会混合两者：用 in-context 快速搭产品，用合成轨迹和 RL 把关键行为训练进模型，再用上下文工程处理生产细节。
\end{knowledgebox}

\subsection{Agent Training 三件套}

技术路线最终都要落到训练和部署，本节拆 Agent training 的三件套。Agent training 的关键环节包括数据合成、强化学习和安全。数据合成生成高质量 trajectory；强化学习依赖清晰 task 和 verifiable reward；安全则需要 sandbox、行为约束和 human-in-the-loop。

\begin{figure}[H]
\centering
\includegraphics[width=0.96\textwidth]{training-ingredients.png}
\caption{Agent Training 三件套：合成轨迹、强化学习和安全约束共同构成训练管线。自制概念图，依据 00:28:29--00:30:57 对谈内容整理。}
\end{figure}

\begin{importantbox}{Agent 数据的范式变化}
传统标注数据常是 input-output；Agent 数据更像 environment + task + action trajectory + reward。模型学的不只是答案，而是如何在环境中行动。
\end{importantbox}

\begin{knowledgebox}{训练环节对照}
\begin{tabularx}{\textwidth}{p{0.22\textwidth}p{0.34\textwidth}X}
\toprule
环节 & 产物 & 主要风险 \\
\midrule
Data Synthesis & 高质量行动轨迹与工具调用示范 & 轨迹不真实，模型学会演戏。 \\
RL & 从环境反馈中优化策略 & reward 设计错误或被 hack。 \\
Safety & 沙盒、权限、人工确认、拒绝策略 & 太松危险，太紧无法完成任务。 \\
Evaluation & 任务成功率、成本、恢复能力 & benchmark 与真实任务不一致。 \\
\bottomrule
\end{tabularx}
\end{knowledgebox}

\subsection{本章小结}

Agent training 的核心是把任务放进环境，让模型通过轨迹、奖励和安全约束学习行动。它不是单纯 SFT，也不是单纯 prompt engineering。

